\subsection{The explicit two-inert continuous class-three obstruction}

The supplied proof of the historical A.196 cites the unavailable
A.181, A.183, and A.195 for its arithmetic nonvanishing.  The following
reconstruction proves that nonvanishing directly in the degree-six
splitting field.  It does not assume the reported degree-nine or
degree-eighteen unit tables, their saturation, or the reported higher
class groups.  Throughout this subsection put
\[
\begin{gathered}
 k=\mathbb Q(\alpha),\qquad \alpha^3=3,\qquad d=11\cdot59=649,\\
 A=d^8\alpha,\qquad U=\alpha^2-2,\qquad B=U^3,\\
 S=\{v:v\mid2\cdot3\cdot11\cdot59\}\cup S_\infty,
 \qquad \Gamma=G_{k,S}.
\end{gathered}
\tag{E.166.1}
\]
Write \(\chi=\chi_3\), \(a=\kappa_3(A)\), and
\(u=\kappa_3(U)\).  Compatible roots may be chosen so that the
literal cocycle for \(B\) is \(b=3u\).  Changing these roots only
changes the prescribed abelian cocycles by coboundaries, which lift
through every quotient used below.

\begin{result}{E.166}{Proved}
{A complete degree-six cup-product certificate}
The map
\[
 a\cup-:H^1(\Gamma,\mathbb F_3)
       \longrightarrow H^2(\Gamma,\mu_3)
\tag{E.166.2}
\]
uses the reduction of \(a\) modulo three and is injective.
Its domain has dimension five.  Separately,
\[
 \overline a\cup\overline u\ne0
       \quad\hbox{in }H^2(\Gamma,\mathbb Z/3(2)).
\tag{E.166.3}
\]
Both assertions are certified by explicit tame residues at the primes
above \(2,11,59\); no calculation in a higher radical field is used.
\end{result}

\begin{proof}
Put \(F=k(\zeta)\), where \(\zeta=\zeta_3\), and let
\(\tau\) fix \(k\) and invert \(\zeta\).  Define
\[
 \delta=1+2\zeta,\qquad
 \beta=\frac{\alpha^2}{\delta},\qquad
 \epsilon=1+\beta-\zeta\beta^2,\qquad
 \eta=\frac{\epsilon}{\tau\epsilon}.
\tag{E.166.4}
\]
Then
\[
 \beta^2=-\alpha,\quad \beta^3=\delta,\quad
 \beta^6=-3,\quad \tau\beta=-\beta,\quad
 N_{F/\mathbb Q(\zeta)}\epsilon=\zeta^2,\quad
 \epsilon\tau\epsilon=U.
\tag{E.166.5}
\]
Thus \(\epsilon\) and \(\eta\) are integral units and
\(\tau\eta=\eta^{-1}\).  These identities follow by reduction
modulo \(\beta^6+3\), or by
\(N(a+b\beta+c\beta^2)=a^3+b^3\delta+c^3\delta^2-3abc\delta\).

We first prove the only class-group fact needed here:
\(\operatorname{Cl}(F)\{3\}=0\).  The field
\(F/\mathbb Q(\zeta)\) is cyclic of degree three, ramified only at
its unique prime above three, and totally ramified there.
The Minkowski bound proves \(h(\mathbb Q(\zeta))=1\).
Let \(H/F\) be its Hilbert three-class field.  It is Galois over
\(\mathbb Q(\zeta)\), with a finite three-group \(P\) as Galois group.
An inertia group at three has order three, because \(H/F\) is
unramified.  Its normal closure is all of \(P\): otherwise the
resulting nontrivial unramified three-group quotient over
\(\mathbb Q(\zeta)\) would have a nontrivial abelian quotient, contrary
to its class number.  A finite three-group normally generated by
one element has cyclic Frattini quotient, hence is cyclic.  Its
inertia subgroup therefore equals \(P\), which has order three.
Consequently \(H=F\).  This argument uses only the Hilbert class
field correspondence, not an asserted numerical class number for
\(F\); see \cite[Chapter V, \S3, especially Example 3.9]{RootMilneCFT}.

Let
\[
 g_2(Z)=Z-1,\quad g_{11}(Z)=Z+2,\quad
 g_{59}(Z)=2+2Z-Z^2,\qquad
 r_q=\frac{g_q(\zeta\alpha)}{g_q(\zeta^2\alpha)}.
\tag{E.166.6}
\]
The exact norms \(N_{k/\mathbb Q}g_q(\alpha)=q\) prove that
\(r_q\) has divisor the difference of the two \(\tau\)-conjugate
primes above \(q\).  At each \(q=2,11,59\), the three primes of
\(F\) have residue field \(\mathbb F_{q^2}\).  Indeed \(q\equiv-1\pmod3\),
and the unique cube root of three in \(\mathbb F_q\) gives the three
roots after adjoining \(\zeta\).  The involution \(\tau\) fixes
one prime and exchanges the other two.  The unique prime above
three is fixed.

Kummer theory, the class-group vanishing just proved, and the fact
that two is invertible modulo three identify
\[
 H^1(\Gamma,\mathbb F_3)
 \simeq
 \bigl(\mathcal O_{F,S_F}^{\times}/
              \mathcal O_{F,S_F}^{\times3}\bigr)^{\tau=-1}.
\tag{E.166.7}
\]
Here the anti eigenspace arises because the chosen identification
\(\mathbb F_3\simeq\mu_3\) over \(F\) changes sign under \(\tau\).
Dirichlet's theorem gives unit rank two for \(F\) and unit rank one
for its \(\tau\)-fixed field \(k\); the minus unit rank is one.
The three-primary roots of unity in \(F\) are exactly \(\mu_3\).
Thus the minus unit classes have dimension two, and the three
minus divisor directions above \(2,11,59\) add three more.
The absence of a three-primary class group makes this valuation
sequence exact modulo three.

The five classes
\[
                   \zeta,\quad\eta,\quad r_2,\quad r_{11},\quad r_{59}
\tag{E.166.8}
\]
are a basis.  The three last classes are independent by their
primitive disjoint divisors.  To check independence of the first
two unit classes, use the two residue maps to \(\mathbb F_{61}\) sending
\(\beta\) respectively to \(3\) and \(19\).  Both send
\(\zeta\) to \(13\); they send \(\eta\) to \(35\) and \(33\).
Writing the cubic residue symbol as \(13^{\lambda}\), the two
rows for \((\zeta,\eta)\) are \((2,1)\) and \((2,0)\).
Their determinant is nonzero modulo three.  The maps are legitimate
on these units: the displayed elements belong to the order generated
by \(\beta,\zeta\), whose index is supported at \(2,3\), and
both denominators in \(\eta\) are units.  Alternatively the
identities in E.166.5 define these two maps directly at primes
above the unramified prime \(61\).

Here is the full tame calculation.  At the prime labeled \((q,j)\),
set \(\alpha=r\zeta^j\) in
\(\mathbb F_{q^2}=\mathbb F_q[\zeta]/(\zeta^2+\zeta+1)\), with
\[
 (q,r)=(2,1),(11,-2),(59,49),\qquad j=0,1,2.
\]
For nonzero local elements \(v,w\), write
\(v=\pi^s v_0\), \(w=\pi^t w_0\).  Our tame-symbol convention is
\[
 (v,w)_3=\left((-1)^{st}v_0^t/w_0^s\right)^{(q^2-1)/3}
          =\zeta^{\lambda(v,w)}.
\tag{E.166.9}
\]
The formula follows from the local Artin description of the Hilbert
symbol, \cite[III, Theorem 4.4 and Remark 4.5, pp.~113--114]{RootMilneCFT},
by evaluating a uniformizer on an unramified cube root of a unit
and using bimultiplicativity and skew symmetry.
Changing the global cup convention reverses all rows and does not
change either conclusion.  Applied to \((A,\zeta),(A,\eta),
(A,r_2),(A,r_{11}),(A,r_{59})\), the rows are
\[
\begin{array}{c|rrrrr|r}
 (q,j)&\zeta&\eta&r_2&r_{11}&r_{59}&(A,U)\\\hline
 (2,0)&0&0&0&0&0&0\\
 (2,1)&0&0&2&0&0&0\\
 (2,2)&0&0&2&0&0&0\\
 (11,0)&1&2&0&1&0&0\\
 (11,1)&1&1&0&0&0&2\\
 (11,2)&1&1&0&0&0&1\\
 (59,0)&2&0&0&1&1&0\\
 (59,1)&2&1&0&1&2&1\\
 (59,2)&2&1&0&1&2&2
\end{array}
\tag{E.166.10}
\]
For clarity, when a factor \(g_q(\zeta^i\alpha)\) vanishes at
the chosen prime, its unit part after dividing by \(q\) is the
inverse product of the other two factors.  This follows from their
exact product \(q\) and removes any need for an unspecified
higher-precision residue lift.  At \(q=11,59\), the unit part
of \(A\) is \((d/q)^8\alpha\) and its valuation is eight.
At two it is simply the unit \(A\).

The five rows \((2,1),(11,0),(11,1),(59,0),(59,1)\) in the
first five columns have integer determinant \(-8\), hence determinant
one modulo three.  Thus every nonzero character in E.166.7 has a
nonzero local cup with \(a\).  Restriction to \(F\) and then to
these local fields proves E.166.2.  The final column gives E.166.3.
All the displayed identities and residue entries are recomputed in
\path{code/check_two_inert_audit.py}; the proof of completeness is
the rank and divisor argument above, not an assumption about the
output of a unit-group algorithm.
\end{proof}

\begin{result}{E.167}{Proved}
{A nonzero supported modulus-nine class and its integral inflation}
Let \(V=\mathbb Z_3[T]/T^3\), with the twisted action
\[
 g f(T)=\chi(g)(1+T)^{a(g)}
                  f\bigl((1+T)^{\chi(g)}-1\bigr).
\tag{E.167.1}
\]
Give \(D_2=\mathbb Z_3[T]/T^2\) the action making the following sequence
exact and equivariant:
\[
 0\longrightarrow D_2\xrightarrow{\ T\ }V
       \xrightarrow{\epsilon}\mathbb Z_3(1)\longrightarrow0.
\tag{E.167.2}
\]
Then the supported connecting class
\(o_S=\partial\kappa_3(B)\) has nonzero reduction in
\(H^2(\Gamma,D_2/9)\), has zero reduction modulo three, and has
exact order three.  Inflation gives a nonzero class
\[
 o\in H^2_{\mathrm{cts}}(G_k,D_2)
\tag{E.167.3}
\]
of the same order.  Its image in
\(H^2(\Gamma,\mathbb Z_3(2))\), and therefore in absolute cohomology,
is zero.
\end{result}

\begin{proof}
Put \(w=1-T/2\).  Direct substitution in E.167.1 gives
\[
 gw=\chi^2w+\chi^2aT,\qquad gT=\chi^3T.
\tag{E.167.4}
\]
Thus
\[
 0\to\mathbb Z_3(3)\xrightarrow{\ T\ }D_2
       \to\mathbb Z_3(2)\to0.
\tag{E.167.5}
\]
Modulo three this is \(0\to\mu_3\to\overline D\to\mathbb F_3\to0\),
where \(\overline D=D_2/3\), with extension cocycle
\(\overline a\).  E.166.2 and its cohomology sequence imply
\[
 H^1(\Gamma,\mu_3)\longrightarrow H^1(\Gamma,\overline D)
                         \quad\hbox{is surjective}.
\tag{E.167.6}
\]

Consider the coefficient Bockstein for
\[
 0\to\overline D\xrightarrow{\times3}D_2/9
                      \longrightarrow\overline D\to0.
\tag{E.167.7}
\]
The last line \(\mu_3\subset\overline D\) lifts to the actual
submodule \(T(\mathbb Z/9)(3)\subset D_2/9\).  Naturality of this
coefficient sequence, together with E.167.6, shows that the entire
Bockstein image from \(H^1(\Gamma,\overline D)\) has zero
projection to \(H^2(\Gamma,\mathbb F_3)\).

Let \(c=\partial_3\kappa_3(U)\in
H^2(\Gamma,\overline D)\), where \(\partial_3\) is the
connecting map of E.167.2 reduced modulo three.  Its bottom
projection is \(\overline a\cup\overline u\), under the
identification of that bottom coefficient with \(\mathbb Z/3(2)\).
It is nonzero by E.166.3.  Consequently \(c\) is not in the
Bockstein image in E.167.7, and
\((\times3)_*c\ne0\).  Apply multiplication by three vertically
to the entire mod-three reduction of E.167.2, mapping it to its
mod-nine reduction.  The maps commute with both \(T\) and
\(\epsilon\).  On the final Tate coefficient the induced map
sends \(\kappa_3(U)\) to \(\kappa_9(U^3)\).  Naturality therefore
gives the exact identity
\[
 (\times3)_*c=\partial_9\kappa_9(U^3)=o_S\bmod9\ne0.
\tag{E.167.8}
\]
Modulo three, \(\kappa_3(B)=0\), so the supported obstruction
is zero.  This proves the asserted first detecting modulus without
any divided-coordinate or higher-field unit calculation.

We next bound its integral order.  For \(F\) as in E.166, let
\(X_0=\operatorname{Spec}\mathcal O_F[1/3]\).
The class-group argument in E.166 and the uniqueness of the prime
above three give
\(H^2_{\acute et}(X_0,\mu_3)=0\): use the Kummer sequence
and the Brauer invariant sequence, whose only possible finite
invariant has sum zero.  The latter sequence is
\cite[II, Proposition 2.1, pp.~163--164]{CohMilneADT}.
For \(j=2,3\), the mod-three sheaf \(\mathbb Z/3(j)\) is isomorphic
to \(\mu_3\) on \(X_0\).  The exact coefficient sequences give
\(H^2_{\acute et}(X_0,\mathbb Z/3^c(j))=0\) for all \(c\).
The group--scheme comparison and finite cohomology used here are
\cite[II, Proposition 2.9, pp.~170--172; I, Corollary 4.15,
p.~62]{CohMilneADT}.  Finite first-cohomology groups remove the
coefficient inverse-limit term.

Removing the primes above \(2,11,59\), the localization sequence
injects the resulting integral second cohomology into
\[
 \bigoplus_{v\mid2\cdot11\cdot59}
 H^1(\mathbb F_v,\mathbb Z_3(j-1))
 =\bigoplus_{v\mid2\cdot11\cdot59}
       \mathbb Z_3/(|\mathbb F_v|^{j-1}-1),\qquad j=2,3.
\tag{E.167.9}
\]
This is the usual codimension-one residue sequence; it follows
also by excision and the tame inertia calculation
\(H^1(I_v,\mathbb Z/3^c(j))=\mathbb Z/3^c(j-1)\), followed by the finite-field
cohomology sequence.  All coefficients have order invertible at
these closed points, and the finite cohomology groups allow the
inverse limit of these exact sequences.  The residue sizes are
\(4,121,3481\); for either exponent \(j-1=1,2\), every summand
in E.167.9 is killed by three.  Restriction--corestriction from
\(F\) to \(k\), of degree two, gives
\[
       3H^2_{\mathrm{cts}}(\Gamma,\mathbb Z_3(j))=0,
                         \qquad j=2,3.
\tag{E.167.10}
\]
The bottom of \(o_S\) is \(a\cup b=3a\cup u=0\).
E.167.5 puts \(o_S\) in the image of
\(H^2(\Gamma,\mathbb Z_3(3))\).  It is killed by three and is
nonzero by E.167.8, so its order is exactly three.

Finally, integral inflation is injective for this particular
module.  Here is a direct continuous-cochain proof.  Let
\(N=\ker(G_k\to\Gamma)\), and suppose a normalized continuous
two-cocycle on \(\Gamma\) becomes \(\delta f\) on \(G_k\).
Its vanishing whenever either argument lies in \(N\) shows that
\(f|_N:N\to D_2\) is an equivariant continuous homomorphism.
The group \(N\) is normally generated by inertia outside \(S\).
At such a place, of residue size \(Q\), wild inertia has no
homomorphism to the three-primary torsion-free module; the only
possible tame image is the image of \(\mathbb Z_3(1)\).
Frobenius equivariance would give
\[
                    (\operatorname{Frob}_v-Q)z=0.
\tag{E.167.11}
\]
The two diagonal eigenvalues on \(D_2\) are \(Q^2,Q^3\).
The determinant \((Q^2-Q)(Q^3-Q)\) is nonzero in \(\mathbb Z_3\),
so E.167.11 forces \(z=0\).  Thus \(f|_N=0\), and the same
normalized equations show that \(f\) descends to a continuous
cochain on the compact quotient \(\Gamma\).  The original
cocycle was already a coboundary.  This proves injectivity and
E.167.3 with its exact order.
\end{proof}

\begin{result}{E.168}{Proved}
{Every finite radical norm exists, and the local norm families are coherent}
Choose compatible roots and put
\[
 K_r=k(A^{1/3^r}),\qquad r\ge0,
 \qquad (A^{1/3^{r+1}})^3=A^{1/3^r}.
\tag{E.168.1}
\]
For every \(r\ge1\), the actual element \(B=U^3\) lies in
\(N_{K_r/k}(K_r^\times)\).  At every place \(v\) of \(k\),
there is a compatible family of completed local unit norms
\[
\begin{gathered}
 b_{r,v}\in\prod_{w\mid v}\widehat{\mathcal O_{K_{r,w}}^\times}_{(3)},\\
 N_{r+1/r}b_{r+1,v}=b_{r,v},\qquad b_{0,v}=B.
\end{gathered}
\tag{E.168.2}
\]
at finite places, with the evident multiplicative local groups at
infinite places.  These statements do not assert a compatible
family of global norm representatives.
\end{result}

\begin{proof}
We first verify the finite local norm conditions explicitly.  At
a finite place not above \(3,11,59\), the element \(A\) is a
unit and \(3\) is invertible.  The algebra defined by
\(Z^{3^r}-A\) is finite etale over the valuation ring, hence is
a product of unramified extensions.  Norms on the units of an
unramified local extension are surjective: finite-field norms are
surjective, and the trace on the residue fields gives successive
lifts on principal units.  As \(B\) is a global unit, it is a
local unit norm at every such place.

At a place over \(11\) or \(59\), the valuation of \(A\) is
eight, prime to three.  The radical extension has degree \(3^r\)
and is totally ramified.  The residue field has size \(q\) or
\(q^2\), where \(q=11,59\); its multiplicative three-primary
part has order one or three.  Cubing therefore kills that part.
The principal units are uniquely \(3^r\)-divisible, by Hensel's
lemma.  Hence the unit \(B=U^3\) is a \(3^r\)-th power in the
base local field.  Such a root, viewed in the degree-\(3^r\)
extension, has norm \(B\).

At the unique place over three, \(v(A)=1\), so
\(Z^{3^r}-A\) is Eisenstein and the local extension has degree
\(3^r\) and is totally ramified.  Its norm is surjective on the
entire base multiplicative group.  Here is the group-theoretic
verification needed to apply the norm limitation theorem.  The
base field \(k_3=\mathbb Q_3(\alpha)\) does not contain \(\mu_3\).
The splitting group embeds in
\((\mathbb Z/3^r)\rtimes(\mathbb Z/3^r)^\times\), with transitive action on
the roots.  Its cyclotomic image contains \(-1\).  A lift of
\(-1\) is an affine involution, and shifting the root labels
by its unique fixed point makes it \((0,-1)\).  For any
\((c,h)\) in the splitting group, its commutator with this
involution is translation by \(\pm2c\).  Since two is a unit,
translation by \(c\) belongs to the group.  Transitivity now
forces the entire translation group \(\mathbb Z/3^r\) to occur.
The involution also places that group in the commutator subgroup.
The root stabilizer maps onto the remaining cyclotomic quotient.
It follows that the radical field has no nontrivial abelian
subextension over \(k_3\).  The local norm limitation theorem
\cite[III, Theorem 3.5, pp.~108--109]{RootMilneCFT} proves norm
surjectivity.  As the extension is totally ramified, a norm
preimage of a unit is a unit.  At real and complex places the
odd degree gives the required norm as well.

We prove the global norm principle in this particular tower;
no higher class-group computation is needed.  Put \(m=3^r\)
and \(L_r=K_r(\mu_m)\).  The field \(k\) has trivial
intersection with every cyclotomic field: a nontrivial intersection
would equal the non-Galois cubic field \(k\) and be abelian over
\(\mathbb Q\).  At a prime over eleven, adjoining \(\mu_m\) is
unramified and the valuation of \(A\) remains eight.  Kummer
theory therefore gives
\[
 G_r=\operatorname{Gal}(L_r/k)
       =(\mathbb Z/m)\rtimes(\mathbb Z/m)^\times,
 \qquad H_r=\operatorname{Gal}(L_r/K_r)=(\mathbb Z/m)^\times.
\tag{E.168.3}
\]
In particular both the translation group and its cyclic complement
are the full displayed groups.

For completeness, the needed cohomology calculation is short.
A generator of the cyclic group \(H_r\) acts on \(\mathbb Z/m\)
by a unit \(h\equiv-1\pmod3\); thus \(h-1\) is invertible.
The homology spectral sequence for E.168.3 has, in total degree
two, only the possible terms
\[
 H_2(H_r,\mathbb Z)=0,\qquad
 H_1(H_r,\mathbb Z/m)=0,\qquad
 H_0(H_r,H_2(\mathbb Z/m,\mathbb Z))=0.
\tag{E.168.4}
\]
The middle equality follows directly from the cyclic resolution:
its degree-one homology is \(\ker(h-1)/N_{H_r}(\mathbb Z/m)\).
Hence \(H_2(G_r,\mathbb Z)=0\), or equivalently
\(H^3(G_r,\mathbb Z)=0\).  Moreover \(G_r^{\mathrm{ab}}=H_r\).
For the norm-one torus \(T_r=R^{(1)}_{K_r/k}\mathbf G_m\),
its character lattice \(J_r\) fits into
\[
 0\to\mathbb Z\to\mathbb Z[G_r/H_r]\to J_r\to0.
\tag{E.168.5}
\]
Shapiro's lemma and the equality of abelianizations show that
\(H^2(G_r,\mathbb Z)\to H^2(H_r,\mathbb Z)\) is an isomorphism.
E.168.4 and the cohomology sequence of E.168.5 give
\[
                         H^2(G_r,J_r)=0.
\tag{E.168.6}
\]

The global duality pairing for tori identifies the dual of
\(\operatorname{Sha}^1(k,T_r)\) with \(\operatorname{Sha}^2(k,J_r)\); see
\cite[I, p.~66, following Corollary 4.21; II, Theorem 4.6 and
Corollary 4.7, pp.~191--192]{CohMilneADT}.
In this case the latter group injects into the finite-group term
E.168.6.  To check that step explicitly, \(J_r\) becomes a
trivial lattice over \(L_r\), so
\(H^1(G_{L_r},J_r)=0\).  Inflation--restriction identifies
the kernel of restriction in degree two with
\(H^2(G_r,J_r)\).  A locally trivial class restricts over
\(L_r\) to a tuple of elements of
\(H^2(G_{L_r},\mathbb Z)=H^1(G_{L_r},\mathbb Q/\mathbb Z)\).
They vanish at every local decomposition group; Chebotarev makes
every such finite-order character zero.  Thus restriction is zero,
and E.168.6 proves \(\operatorname{Sha}^1(k,T_r)=0\).
Hilbert 90 identifies \(H^1(k,T_r)=k^\times/N_{K_r/k}K_r^\times\).
The local norms already verified consequently give an actual
global norm preimage of \(B\) at every finite height.
The same conclusion is a special case of
\cite[Theorems 1.2 and 3.2, pp.~2 and 6]{LitHoshiYamasaki2025Metacyclic};
that work is a preprint, and the group calculation and the
established duality argument above supply the proof used here.

Finally fix a finite place \(v\).  The three-primary completions of
the local unit groups at each height are compact.  The local
norm preimages proved above give a solution to every finite
initial compatibility problem: take a preimage at its largest
height and norm it down.  The conditions are closed.  Compactness
of the product gives E.168.2.  This compactness argument is local;
it is not applied to the unrestricted global field groups.  At an
infinite place the odd-primary completed multiplicative group is
zero, so the compatible completed family is the trivial one.
\end{proof}

\begin{result}{E.169}{Proved}
{The exact A.196 continuous class-three phenomenon}
The actual pro-three tripod quotient over \(k\) with geometric
group
\[
 \mathcal Q=\bigl(\mathbb Z_3[T]/T^3\bigr)y\rtimes\mathbb Z_3x,
                   \qquad xvx^{-1}=(1+T)v
\tag{E.169.1}
\]
has lower-central class three and abelian derived subgroup.
The prescribed pair \((\kappa_3(A),\kappa_3(B))\) has a global
continuous class-two lift and a continuous local
\(\mathcal Q\)-lift at every place.  It has no global
continuous \(\mathcal Q\)-lift.  The obstruction is the
specific order-three class \(o\) of E.167.

For every \(c\ge1\), the same pair does have a global lift to
\[
 \mathcal Q_c=\bigl((\mathbb Z/3^c)[T]/T^3\bigr)\rtimes\mathbb Z/3^c,
               \qquad |\mathcal Q_c|=3^{4c}.
\tag{E.169.2}
\]
These quotients have exact class three and are cofinal in
\(\mathcal Q\).  Their global lift classes have no compatible
inverse family.  Every fixed finite ramification set \(S'\supset S\)
fails at some finite coefficient level and at every higher level;
for \(S'=S\), the first failing level is exactly \(c=2\).
All finite absolute coefficient images of \(o\) are zero, so
it is a nonzero order-three element in the coefficient
\(\varprojlim^1 H^1\) subgroup.

This is a statement over \(\mathbb Q(\sqrt[3]{3})\), for a geometric
pro-three quotient.  Its prescribed local abelian pair at three
is not ordinary or cuspidal modulo nine, as proved in E.43.
It supplies neither an ordinary completed pair over \(\mathbb Q\)
nor a section of its full pronilpotent quotient.
\end{result}

\begin{proof}
We retain the actual tangential arithmetic action.  For the free
geometric pro-three group \(\Pi=\langle x,y\rangle\), put
\(N=\ker(\Pi\to\mathbb Z_3)\), with \(x\mapsto1,y\mapsto0\).
The pro-three Schreier calculation gives
\[
 \Pi/[N,N]=\mathbb Z_3[[T]]y\rtimes\mathbb Z_3x,
                         \qquad xvx^{-1}=(1+T)v.
\tag{E.169.3}
\]
One direct verification is to abelianize the normal subgroup in
each cyclic cover: its conjugates \(x^i yx^{-i}\) give the
permutation basis, and the compatible trace maps give the
completed group algebra.  The integral tangential normalization
has \(g(x)=x^\chi\), \(g(y)=f_g^{-1}y^\chi f_g\), with
\(f_g\in\Pi'\subset N\).  The latter conjugation disappears
in \(N^{\mathrm{ab}}\).  Thus the induced action is precisely
\[
 g(v,t)=\bigl(\chi\,v((1+T)^\chi-1),\chi t\bigr).
\tag{E.169.4}
\]
The ideal \((T^3)\) is stable.  Quotienting E.169.3 by it gives
the actual arithmetic quotient E.169.1, rather than a replacement
of its Galois action.  The tangential action used here is the
same one fixed in E.126--E.127.

Commuting with \(x\) multiplies the normal abelian module by
\(T\).  Hence
\(\Gamma_2\mathcal Q=TR\), \(\Gamma_3\mathcal Q=T^2R\ne0\),
\(\Gamma_4\mathcal Q=0\), with \(R=\mathbb Z_3[T]/T^3\).
The same identities hold modulo \(3^c\).
For comparison with the full pro-three class-two quotient, write
\(v=b_0+c_0T+e_0T^2\) and \(z_0=e_0+c_0/2\).
The group is the matrix group
\[
 (a_0,b_0,c_0,z_0)\longmapsto
 \begin{pmatrix}
 1&a_0&a_0^2/2&z_0\\
 0&1&a_0&c_0\\
 0&0&1&b_0\\
 0&0&0&1
 \end{pmatrix},
\tag{E.169.5}
\]
with arithmetic weights \(\chi,\chi,\chi^2,\chi^3\).
Discarding \(z_0\) is the actual Heisenberg quotient
\(\Pi/\Gamma_3\Pi\).  A section with first cocycle \(a\)
is exactly a cocycle in the twisted module E.167.1; its second
abelian coordinate is its augmentation.  Thus its obstruction
is \(\partial b=o\).  Abelian coboundaries can be removed by
lifting zero-cochains through the \(x\)-coordinate and through
augmentation.  The obstruction consequently applies both to
literal cocycles and to their geometric-conjugacy classes.
E.167 proves both the global class-two lift and the failure of
all global continuous class-three lifts.

At an infinite place, odd-primary first cohomology is zero and the
prescribed abelian Kummer classes are zero; choose the trivial
section class.  For a finite place, use the compatible family
E.168.2 and Shapiro's lemma on each radical root set.  The normal modules at finite
root height are
\[
 P_r=\mathbb Z_3[X]/(X^{3^r}-1),\qquad
 g f(X)=\chi(g)X^{a(g)}f(X^{\chi(g)}).
\tag{E.169.6}
\]
These are the induced Tate modules of \(K_r/k\); augmentation
on their first cohomology is the actual field norm.  Norm
compatibility in E.168.2 is therefore compatibility of the
local Shapiro classes.  The trace transition maps of these
modules are surjective.  Their local invariants are
\(H^0(K_r\otimes k_v,\mathbb Z_3(1))=0\) at finite places, since
a finite local field has only finitely many three-power roots
of unity.
The degree-one cochain inverse-limit sequence now lifts the
compatible classes to a continuous cocycle in
\(\varprojlim P_r=\mathbb Z_3[[T]]\), with augmentation \(b\).
Its quotient by \(T^3\) is the claimed local section.
Equivalently one may form the inverse limit of its finite local
cocycle fibres; those fibres are finite and have the same closed
compatibility conditions.

For the finite global lift, put \(m=3^c\).  Since both
\(m\) and \(\binom m2\) are divisible by \(3^c\),
\[
                   X^{3^c}-1\in(3^c,T^3),\qquad X=1+T.
\tag{E.169.7}
\]
The actual norm preimage of \(B\) from \(K_c\) supplied by
E.168 has a Shapiro cocycle in \(P_c\).  The map
\(X\mapsto1+T\), with coefficient and \(x\)-coordinate
reduction modulo \(3^c\), takes its section to
\(\mathcal Q_c\).  This proves every claimed finite absolute
lift and, by the connecting sequence, the vanishing of every
absolute image \(o\bmod3^c\).

The four matrix coordinates in E.169.5 give the order in
E.169.2 and identify their kernels as a neighborhood basis in
\(\mathcal Q\).  A compatible sequence of literal finite
sections would have a continuous inverse limit, contradicting
E.167.  Compatible geometric-conjugacy classes also give such
a sequence: inductively lift a conjugating element through
\(\mathcal Q_{c+1}\to\mathcal Q_c\) to make representatives
compatible.  Their abelian classes agree with \((a,b)\) at
every level.  The correcting abelian zero-cochains are unique,
because \(H^0(k,\mu_{3^c})=0\), and hence compatible.
The limit would have the prescribed integral abelian class.

If a finite set \(S'\) is fixed, the action and prescribed
frame factor through \(G_{k,S'}\).  For each \(c\) its literal
cocycle lifts form a finite set: adjoining the finite image
of the action turns a cocycle into a homomorphism to a finite
semidirect product, and Hermite finiteness bounds the number
of extensions of bounded degree ramified only in \(S'\).
Nonempty lift sets at unbounded levels would, by reduction,
solve every finite compatibility problem.  Compactness of
the product of these finite sets would contradict the
nonexistence just proved.  Thus some level and all higher
levels fail.  For \(S'=S\), \(b\bmod3=0\) gives the
level-one section \(x^{a\bmod3}\), whereas E.167.8 excludes
level two.  This is the exact first failure.

Finally the surjective finite coefficient maps give the
continuous-cochain exact sequence
\[
 0\longrightarrow\varprojlim\nolimits_c^1
                  H^1(G_k,D_2/3^c)
 \longrightarrow H^2_{\mathrm{cts}}(G_k,D_2)
 \longrightarrow\varprojlim_c H^2(G_k,D_2/3^c)
 \longrightarrow0.
\tag{E.169.8}
\]
It follows directly by applying cohomology to the usual
product-complex map \(1-\mathrm{shift}\); continuous finite-valued
cochains lift through each surjective coefficient map.
The nonzero class \(o\), now explicitly constructed and proved
to have order three, lies in the first term.  E.43 supplies
the stated local ordinary and cuspidal exclusion, so no
ordinary-point hypothesis is inferred from these local
quotient sections.
\end{proof}
